\section{Review Questions}

\begin{enumerate}
\item Explain a linear system as a collection of simultaneous conditions. What does it mean for an ordered pair to belong to the solution set?

\item Interpret a $2\times2$ linear system geometrically. Describe the configurations corresponding to one solution, no solution, and infinitely many solutions.

\item Write a given pair of linear equations in matrix form $Ax=b$ and as an augmented matrix. Explain what information belongs to each block.

\item Solve one simple $2\times2$ system by Gaussian elimination. At every row operation, state why the new system has the same solution set as the previous one.

\item Explain the roles of pivots and echelon form. How do pivots identify independent conditions and determined variables?

\item What algebraic pattern signals inconsistency? Give a $2\times2$ example and connect the contradictory row with the geometry of parallel lines.

\item How does a free variable arise in a dependent $2\times2$ system? Parameterise the solutions of one example and explain geometrically what the parameter measures.

\item State Cramer's rule for a $2\times2$ system. Why must the coefficient determinant be nonzero, and why does the rule not classify singular systems by itself?

\item Explain the inverse-matrix method $x=A^{-1}b$. Compare its meaning and practical role with Gaussian elimination and Cramer's rule.

\item Classify the parameter-dependent system
\[
\begin{cases}
x+y=3,\\
2x+ay=b.
\end{cases}
\]
Identify the generic and exceptional parameter cases, and explain why determinant information must sometimes be followed by a compatibility check.
\end{enumerate}
